\section{Discussion and Limitations}
\label{sec:discussion}

The results identify two reasons to consider the proposed Fuzzy PI law
for bittide: a provable reduction in transient buffer demand for a
specified reference execution, and independently tested synchronization
benefits in defined delivery regimes. The connection between these
results is the explicit nonlinear correction mechanism, rather than
an assumption that fuzzy inference improves every performance metric.
This section explains how the benefits should inform controller
selection and which extensions require further evidence.

\subsection{What nonlinear response shaping adds}
\label{subsec:discussion_mechanism}

The controller offers a local way to change the relationship between
occupancy error, its change between measurements and the next frequency
correction. Proposition~\ref{prop:fpi_central} shows that the normalized
min-activation surface strengthens its calibrated linear counterpart
for aligned input signs and attenuates its magnitude for opposing signs
within the stated central region. A uniform increase in PI gains would
scale both cases uniformly. The Fuzzy law therefore supplies
a different response shape at calibrated gain coordinates, while
retaining a concrete connection to the native PI implementation.

This interpretation concerns a correction increment, not a complete
closed-loop performance ordering. In particular, aligned signs do not
always mean that absolute occupancy error is increasing; an error can
cross zero while decreasing in magnitude. The plant, measurement
history and accumulated correction determine the resulting trajectory.
Equation~\eqref{eq:fpd_accounting} makes the memory contribution explicit:
past nonlinear increments may remain in the accumulator even when
later inputs have zero nonlinear residual. This provides a mechanism
to investigate transient shaping and persistent offsets, but a
same-input decomposition alone does not establish their causal
contribution to differences between independently evolving closed loops.

The comparisons answer complementary questions. The matched-gain
reference proof and startup experiments hold the calibrated coordinates
fixed. They test a prescribed change in the controller law. The
equal-budget study instead asks whether each controller family can
produce a useful selected design with the same finite search budget.
Its lower held-out RMS supports the selected Fuzzy scheme, including
its chosen gains, rather than an inference-only effect at identical
gains. Keeping these questions separate makes the empirical benefit
more meaningful than a comparison against a single arbitrary PI setting.

\subsection{Engineering meaning of the buffer-headroom benefit}
\label{subsec:discussion_headroom}

For a finite receiving buffer, the largest transient deviation matters
even when the eventual operating point is admissible.
Theorem~\ref{thm:btd_headroom} compares a certified upper bound for the
Fuzzy excursion with an attained PI peak, not just two conservative
upper bounds. For its prescribed initialization, the resulting
improvement in minimum headroom exceeds 6.0887 fluid frames. Both
controllers also satisfy the reference model's positive-frequency,
correction and convergence conditions. The benefit thus concerns
transient resource use while those conditions remain satisfied.

The matched-gain actual tests provide a separate engineering observation.
For example, at 1600 ppm, Fuzzy reduces the measured excursion from
94 to 72 frames without correction saturation or gate intervention
(Table~\ref{tab:eval_startup_safety}). This is consistent with a useful
headroom effect after introducing delayed local execution and integer
buffer events, although the mathematical reference certificate does
not cover that implementation. At the larger stress scale, fewer
saturated updates and smaller excursions remain finite-run findings,
not a proof that saturation is harmless or that future overflow is
impossible.

These results measure reserve at fixed capacity. They do not establish
that a smaller hardware buffer can simply be substituted. Capacity
appears in the feedback normalization
$\gamma_i=100/(d_iC)$, where $d_i$ is the number of receiving neighbours.
Changing $C$ or the degree changes the token-domain controller map
unless the gains or normalization are adjusted consistently. Any
capacity-reduction claim would therefore require a new model-consistent
comparison and a fresh safety check. Similarly, a failed sufficient
condition is inconclusive: the observed-safe but uncertified reference cases
illustrate conservatism, not an enlarged general safe set.

\subsection{Selecting a controller for the operating regime}
\label{subsec:discussion_selection}

The independent tests demonstrate a practical benefit under the stated
selection policy. With identical evaluation budgets, frozen gains and
shared untouched profiles, the random-priority Fuzzy representative
reduces disagreement RMS and mean command TV in both original random
regimes. The burst-priority representative also substantially reduces
IAE in the larger held-out burst. The distinct selected representatives
show that the desired task matters: prioritizing random disagreement
and prioritizing a burst transient need not choose the same parameters.
The experiments do not switch between these representatives online.

Several measures are needed to interpret the benefit. Relative
frequency agreement concerns differences between nodes; it does not
require the common frequency to equal one or to return to its clean
counterfactual. The larger common-frequency shift accompanying the
selected Fuzzy RMS improvement is therefore an engineering cost to
evaluate against the application, not an inconsistency in the RMS
measurement. The controller has no absolute-frequency reference loop.
Likewise, smaller buffer peaks need not coincide with the lowest RMS,
and command TV is an activity metric rather than a measured energy
cost. A design choice should specify the acceptable headroom,
frequency-shift and correction constraints before selecting among the
reported tradeoffs.

The slower-delivery study further shows why timing conditions belong
in that choice. Increasing a constant path-delay plateau does not
suspend arrivals for the whole plateau: its onset changes in-flight
inventory, and the first deficit can form before changed feedback is
applied. The ramp profiles allow regulation during the formation of
the imbalance, whereas the random profile grids change its temporal
structure. Neither a longer hold nor slower variation alone establishes
an RMS advantage. These observations motivate matching the controller
to a declared disturbance regime; they do not identify an optimal
polling period, since polling was not independently swept in the new
performance tests.

\subsection{Meaning of safety, convergence and logical synchrony}
\label{subsec:discussion_scope}

The certificate applies to the synchronous zero-delay fluid model with
its compatible targets, initialization and invariant input region.
Its sampled recurrence is the exact endpoint map of that continuous
zero-order-hold reference. It is not a replacement discrete firing-count
plant, nor an exact model of the asynchronous controller scheduling
and integer delivery events in actual Callisto. In particular, the
complete 25-rule controller remains executable outside the central
input region, while the regional formula and the numerical invariant
certificate cannot be applied to arbitrary such trajectories.

The actual tests establish capacity compliance, positive frequency and
progress over specified finite records. Persistent random input is
evaluated by windowed disagreement and occupancy statistics rather
than claimed exact convergence to a fixed equilibrium. The startup
records also have a shorter declared horizon than the disturbance
tests. None of these observations alone supplies an all-time safety
or convergence proof for delayed execution. Similarly, the reference
frequency-spread result requires a separately justified engineering
threshold; logical synchrony does not supply a universal ppm limit.

Logical synchrony is the computational requirement motivating the
controller, while physical syntony, buffer centering and positive
frequency describe different properties of its implementation.
Preserving capacity and frame order, together with continued data
progress, supports the normal execution assumed here. Demonstrating
functional logical synchrony for an application would additionally
require checking the logical frame relationships and payload execution
against the corresponding logical model. These experiments do not
perform that refinement check or measure application-level useful
throughput. A capacity guard that prevents underflow by stopping data
progress must also be distinguished from successful unstalled control.

\subsection{External validity and next steps}
\label{subsec:discussion_future}

The independent comparison is deliberately bounded: a three-node
triangle, a fixed two-dimensional gain grid, fixed Fuzzy input scaling
$K_e$, and one training seed per random regime. The untouched seeds
quantify variation within those declared input families; they do not
turn the search into global controller optimization or establish
performance on larger, heterogeneous or changing networks. The
transport model preserves order and permits batch release without a
finite serialization rate. Hardware oscillator limits, serialized
links, packet loss, node dropout and controller execution cost require
separate evaluation. The current observer and supervisory switching
mechanisms are disabled, so they contribute no fault-recovery claim.

The most direct theoretical extension is to account explicitly for
local sampling, bounded delay and quantization when constructing
invariant buffer and input regions. Such a result should preserve
absolute-frequency and actuator constraints as well as disagreement
bounds. Empirically, broader topologies and delivery models would
test whether the observed benefits persist beyond the present
triangle. If absolute-rate requirements or sustained saturation are
important, a common-frequency constraint or anti-windup mechanism
could be studied as a separate controller extension, followed by
fresh analysis and independent tests. These steps extend the present
mechanism and evidence without changing the scope of its proved or
measured benefits.


% Previous placeholder, preserved before the completed section was integrated.
% \section{Discussion and Limitations}
% \label{sec:discussion}
% \noindent\textbf{Draft status: to be written.}
% The full discussion remains pending. The certificate in
% Section~\ref{sec:btd_reference} applies to its specified synchronous,
% zero-delay fluid reference and prescribed initialization. It does not
% certify asynchronous, delayed, quantized Callisto execution or establish
% universal superiority over PI.
